\section{Scope and relation to the THUMS series}

\subsection{Why a new series rather than a new version}

P2H2P/THUMS is frozen at v0.16 and remains runnable. Nothing in BOREAS edits
it; the build script hashes the THUMS sources against their recorded stamp and
refuses to proceed if they have moved, so the promise is enforced rather than
asserted. The separation is deliberate. The geometry change is not a refinement
of the hairpin model but a different machine, and the two should be comparable
to one another rather than one silently overwriting the other. In particular,
the \SI{92}{\kilo\watt} parasitic quantified in Section~\ref{sec:hairpin} is a
\emph{result} about the previous geometry, and results require the previous
geometry to remain executable.

\subsection{What is inherited unchanged}

The following are carried over from THUMS v0.16 without modification, and are
documented in the companion model report rather than repeated here:

\begin{itemize}[nosep]
\item the enthalpy formulation (``Formulation C''), in which the stored state
      is $\Ep$, the enthalpy per unit length of tube measured from a
      fully-solid-at-$\Tm$ datum, with the three-branch constitutive curve
      giving subcooled solid, two-phase and superheated liquid;
\item the cyclic steady state (CSS) solver, which marches charge and discharge
      alternately until the end-of-cycle state repeats to tolerance;
\item the quasi-steady melt-layer conduction model and the direction-dependent
      choice of conduction path;
\item the double-stage ORC and the high-temperature heat pump, including
      isentropic machine efficiencies applied inside the cycles;
\item the component-wise exergy audit, the Gouy--Stodola accounting, and the
      treatment of the finite geothermal resource;
\item the single ORC correction pass that resolves the coupling between the
      evaporating temperature and the realised discharge outlet.
\end{itemize}

\subsection{What is new}

Four things, in the order they are developed below: the downcomer tube
geometry (Section~\ref{sec:geometry}); the reinterpretation of the axial
coordinate as a depth, with the cascade orientation that follows
(Section~\ref{sec:coordinate}); the conducting formation
(Section~\ref{sec:formation}); and the depth-dependent field shielding
(Section~\ref{sec:field}).

